\documentclass[11pt]{article}

\usepackage[margin=1in]{geometry}
\usepackage[T1]{fontenc}
\usepackage[utf8]{inputenc}
\usepackage{lmodern}
\usepackage{amsmath,amssymb,amsthm}
\usepackage{hyperref}
\usepackage{xcolor}
\usepackage{listings}
\usepackage{multicol}

\definecolor{codebg}{RGB}{248,248,248}
\definecolor{codeframe}{RGB}{210,210,210}
\definecolor{linkblue}{RGB}{0,72,130}

\hypersetup{
  colorlinks=true,
  linkcolor=linkblue,
  citecolor=linkblue,
  urlcolor=linkblue
}

\lstdefinelanguage{Lean}{
  keywords={
    abbrev,by,def,deriving,end,exact,exists,false,forall,have,import,
    inductive,namespace,not,rcases,structure,theorem,true,where,with
  },
  sensitive=true,
  comment=[l]{--},
  morecomment=[s]{/-}{-/},
  morestring=[b]"
}

\lstset{
  basicstyle=\ttfamily\scriptsize,
  backgroundcolor=\color{codebg},
  frame=single,
  rulecolor=\color{codeframe},
  breaklines=true,
  columns=fullflexible,
  keepspaces=true,
  tabsize=2
}

\newtheorem{definition}{Definition}
\newtheorem{invariant}{Invariant}
\newtheorem{theorem}{Theorem}
\newtheorem{lemma}{Lemma}

\title{\textbf{ic-memory}\\
Stable-Memory Allocation Governance for Upgrade-Safe Internet Computer Canisters}
\author{Adam Powell}
\date{\today}

\begin{document}
\maketitle

\begin{abstract}
Internet Computer canisters often store durable state through
\texttt{ic-stable-structures}. Multiple stable stores are commonly multiplexed
through a \texttt{MemoryManager}, where each logical store is assigned a physical
stable-memory slot. Across upgrades, an accidental reassignment can make a
canister open one store's bytes as another store's data. \texttt{ic-memory}
addresses this class of failure by maintaining a durable allocation ledger that
binds stable logical keys to physical \texttt{MemoryManager} IDs and rejects
layout drift before stable-memory handles are opened. This paper presents the
problem, protocol, core invariants, and a compact Lean mechanization of the
allocation-safety argument.
\end{abstract}

\begin{multicols}{2}
\raggedright

\section{Problem Statement}

Stable memory is durable across Internet Computer canister upgrades. That
durability is useful only when future binaries interpret the same durable bytes
under the same logical identity. Consider a canister version that ships with
\[
  \begin{aligned}
  \texttt{app.users.v1} &\mapsto 100,\\
  \texttt{app.orders.v1} &\mapsto 101.
  \end{aligned}
\]
If a later binary accidentally ships with the assignments swapped, the program
may compile and install while opening orders data as users data. The failure is
not a memory-access violation; it is an allocation-ABI violation.

\texttt{ic-memory} treats the mapping
\[
  \text{stable key} \longrightarrow \text{physical allocation slot}
\]
as persistent protocol state. A future binary may reopen the same key on the
same slot, or it may introduce a new key. It may not move an existing key to a
new slot, and it may not reuse an existing slot for another active key.

\section{Scope}

The crate is allocation-governance infrastructure. It does not implement stable
collections, prove application schema compatibility, authorize controllers, or
perform data migrations. Its narrower responsibility is to validate declarations
against historical allocation facts before opening memory.

The native substrate is the \texttt{MemoryManager} from
\texttt{ic-stable-structures}. Its usable ID domain is
\[
  S = \{0,1,\ldots,254\},
\]
with ID \(255\) reserved by \texttt{ic-stable-structures} as the unallocated
sentinel. In the default runtime, \texttt{ic-memory} reserves IDs
\texttt{0..=9} for governance and uses ID \(0\) for the allocation ledger.
Lower-level bootstrap
owners that bypass the default runtime must enforce equivalent policy
themselves.

\section{Substrate Correspondence}

The allocation slot model is intentionally aligned with the current
\texttt{ic-stable-structures} \texttt{MemoryManager} implementation. The
upstream manager defines \texttt{MemoryId} as a one-byte value, uses
\texttt{255} as the unallocated-bucket marker, and stores each bucket owner as
one byte in the manager metadata. Therefore \texttt{ic-memory}'s current slot
descriptor is not an arbitrary namespace: it is the exact durable ID domain that
the underlying manager can distinguish.

The same implementation reserves the first stable-memory page for manager
metadata. Its V1 header stores the magic value \texttt{MGR}, a layout version,
the allocated-bucket count, the bucket size, 32 reserved bytes, and one
page-count entry for each managed memory \(0,\ldots,254\). The bucket-owner
table follows the header. With the default bucket size of 128 Wasm pages and
32,768 managed buckets, the manager can address 256 GiB of bucket-backed stable
memory through this layout.

\texttt{ic-memory}'s default ledger anchor also relies on
\texttt{ic-stable-structures::Cell}. The relevant Cell V1 envelope begins with
magic \texttt{SCL}, a one-byte layout version, and a four-byte value length
before the encoded value. \texttt{ic-memory} preflights that envelope before
opening the ledger Cell so corrupt ledger storage is classified as a bootstrap
error instead of escaping as a decode panic.

\section{Protocol Overview}

Each binary contributes a declaration snapshot \(D\), containing pairs
\[
  (k,s) \in K \times S
\]
where \(k\) is a stable key such as \texttt{app.users.v1}, and \(s\) is a
physical \texttt{MemoryManager} ID. The runtime then recovers the committed
allocation ledger, collects current declarations, validates them against history
and policy, commits the next generation, and publishes a capability that can
open validated slots.

The default runtime performs that sequence before publishing its open authority.
The lower-level Rust APIs expose the pieces separately for framework owners, so
manual integrations must preserve the same order.

The ordering is the central safety boundary:
\[
  \boxed{
    \begin{gathered}
    \text{recover} \rightarrow \text{validate}\\
    \rightarrow \text{commit} \rightarrow \text{open}
    \end{gathered}
  }
\]
Opening stable-memory handles before this boundary defeats the protocol.

\section{State Model}

Let \(K\) be the set of stable keys and \(S\) the set of usable physical slots.
A ledger is a finite sequence of records
\[
  L = [r_1,\ldots,r_n].
\]
Each record has
\[
  r = (k, s, state, first, last, retiredAt)
\]
where \(k \in K\), \(s \in S\), and
\[
  state \in \{\mathsf{Reserved}, \mathsf{Active}, \mathsf{Retired}\}.
\]

\begin{definition}[Active binding]
A stable key \(k\) is active at slot \(s\) in ledger \(L\), written
\(\mathsf{ActiveAt}(L,k,s)\), when
\[
  \exists r \in L.\; r.key = k \land r.slot = s
  \land r.state = \mathsf{Active}.
\]
\end{definition}

\begin{definition}[Retired binding]
A stable key \(k\) is retired at slot \(s\), written
\(\mathsf{RetiredAt}(L,k,s)\), when
\[
  \exists r \in L.\; r.key = k \land r.slot = s
  \land r.state = \mathsf{Retired}.
\]
\end{definition}

\section{Allocation Invariants}

\begin{invariant}[No stable-key movement]
For all records \(r_1,r_2 \in L\),
\[
  r_1.key = r_2.key \Rightarrow r_1.slot = r_2.slot.
\]
Once a stable key is assigned to a slot, that key cannot later name another
slot.
\end{invariant}

\begin{invariant}[No physical-slot reuse]
For all records \(r_1,r_2 \in L\),
\[
  r_1.slot = r_2.slot \Rightarrow r_1.key = r_2.key.
\]
Once a physical slot is assigned to a stable key, the slot cannot be reused for
a different stable key.
\end{invariant}

\begin{invariant}[Retirement is a tombstone]
For all \(k \in K\) and \(s \in S\),
\[
  \mathsf{RetiredAt}(L,k,s) \Rightarrow
  \neg \mathsf{ActiveAt}(L,k,s).
\]
Retirement is not a free-list operation. It preserves historical ABI facts for
rollback safety and diagnostics.
\end{invariant}

\begin{invariant}[Post-commit open authority]
In the abstract model, a memory handle is opened through a post-commit authority
tied to a committed ledger. If that authority permits opening \(k\) at \(s\),
the committed ledger contains an active record for \(k\) at \(s\):
\[
  \mathsf{MayOpen}(A,k,s) \Rightarrow \mathsf{ActiveAt}(L,k,s).
\]
For the Rust default runtime, this corresponds to publishing validated
allocations only after the staged ledger generation has been committed. Manual
integrations carry this as an ordering obligation.
\end{invariant}

\begin{invariant}[Generation monotonicity]
Successful staging advances the committed ledger generation by exactly one in
the abstract model:
\[
  g' = g + 1 \quad \text{and therefore} \quad g < g'.
\]
The Rust implementation refines this with explicit checks for stale validated
generations, declaration-count bounds, and \texttt{u64} overflow.
\end{invariant}

\section{Proof Sketch}

\begin{theorem}[Stable key has a unique slot]
Assume a ledger \(L\) satisfies no stable-key movement. If records \(r_1,r_2\)
are both in \(L\) and \(r_1.key = r_2.key\), then \(r_1.slot = r_2.slot\).
\end{theorem}

\begin{proof}
This is direct application of the no stable-key movement invariant. The
interesting implementation obligation is therefore not the theorem itself, but
ensuring every recovery, validation, reservation, staging, and commit path
preserves the invariant before producing authority.
\end{proof}

\begin{theorem}[Physical slot has a unique stable key]
Assume a ledger \(L\) satisfies no physical-slot reuse. If records \(r_1,r_2\)
are both in \(L\) and \(r_1.slot = r_2.slot\), then \(r_1.key = r_2.key\).
\end{theorem}

\begin{proof}
This is direct application of the no physical-slot reuse invariant. The result
prevents the swapped-ID failure: two distinct keys cannot both own the same
stable-memory slot.
\end{proof}

\begin{theorem}[Open authority is ledger-backed]
Let \(A\) be a post-commit open authority tied to committed ledger \(L\). If
\(A\) permits opening stable key \(k\) at slot \(s\), then \(L\) contains an
active allocation record for \(k\) at \(s\).
\end{theorem}

\begin{proof}
By model construction, a post-commit open authority contains only declarations
whose matching allocation records are active in the committed ledger. Opening
resolves a stable key to one of those declarations, so the corresponding active
ledger record exists.
\end{proof}

\begin{theorem}[Generation advancement is monotone]
If staging succeeds at generation \(g\), the staged generation is \(g+1\), and
\[
  g < g + 1.
\]
\end{theorem}

\begin{proof}
The generation counter advances by successor in the model. The Rust
implementation has additional failure cases for stale validated generations,
declaration-count limits, and \texttt{u64} overflow; this theorem should not be
read as a verification of those checks.
\end{proof}

\section{Lean Mechanization}

The executable model is in \path{whitepaper/lean/IcMemory.lean}. It proves the
core protocol facts over a compact abstraction of the ledger. The model avoids
Rust-specific implementation details and instead captures the authority
boundary that the Rust code must maintain. In particular, the Lean
\texttt{OpenAuthority} type is a post-commit model object, not a field carried
inside Rust's \texttt{ValidatedAllocations} struct.

\begin{lstlisting}[language=Lean,caption={Core allocation records and ledger invariants.}]
abbrev StableKey := Nat
abbrev Slot := Fin 255

inductive AllocationState where
  | reserved
  | active
  | retired

structure Record where
  key : StableKey
  slot : Slot
  state : AllocationState
  firstGeneration : Nat
  lastSeenGeneration : Nat
  retiredGeneration : Option Nat

def NoKeyMovement (records : List Record) : Prop :=
  forall r1, r1 in records ->
  forall r2, r2 in records ->
    r1.key = r2.key -> r1.slot = r2.slot
\end{lstlisting}

\begin{lstlisting}[language=Lean,caption={The checked no-movement theorem.}]
theorem stable_key_has_unique_slot
    (ledger : SafeLedger)
    {r1 r2 : Record}
    (h1 : r1 in ledger.records)
    (h2 : r2 in ledger.records)
    (sameKey : r1.key = r2.key) :
    r1.slot = r2.slot :=
  ledger.noKeyMovement r1 h1 r2 h2 sameKey
\end{lstlisting}

\begin{lstlisting}[language=Lean,caption={Post-commit opening is backed by the ledger.}]
theorem authority_open_is_backed_by_active_ledger_record
    {ledger : SafeLedger}
    (authority : OpenAuthority ledger)
    {key : StableKey}
    {slot : Slot}
    (openable : MayOpen authority key slot) :
    ActiveAt ledger.records key slot := by
  rcases openable with <declaration, inAuthority, sameKey, sameSlot>
  have backed := authority.sound declaration inAuthority
  ...
\end{lstlisting}

The TeX listings use ASCII approximations for readability. The source file uses
Lean syntax accepted by \texttt{lake build}.

\section{Durable Commit Protocol}

The allocation ledger is stored under the default native anchor
\[
  \begin{gathered}
  \texttt{MemoryManager ID 0}\\
  \rightarrow \texttt{Cell<StableCellLedgerRecord>}\\
  \rightarrow \texttt{LedgerCommitStore}.
  \end{gathered}
\]
\texttt{LedgerCommitStore} uses dual protected commit slots. Each slot contains
a generation number, marker, checksum, and encoded ledger payload. Recovery
chooses the highest-generation valid slot. A corrupt newer slot cannot override
an older valid slot, and ambiguous equal-generation divergent slots fail closed.

The logical payload is wrapped in a small envelope before ledger decode:
\[
  \begin{gathered}
  \texttt{ICMEMLED}
  \parallel \texttt{envelopeVersion}\\
  \parallel \texttt{ledgerSchemaVersion}
  \parallel \texttt{physicalFormatId}\\
  \parallel \texttt{payloadLength}
  \parallel
  \texttt{CBOR(AllocationLedger)}.
  \end{gathered}
\]
This ordering matters: version routing happens before deserializing the ledger
DTO. The decoded DTO still is not authority; it becomes useful for allocation
authority only after compatibility checks, committed-integrity validation, and
\texttt{RecoveredLedger} construction succeed.

\section{Operational Guidance}

Use \texttt{ic-memory} when a canister has multiple stable stores, generated
stores, framework-owned stores, or plugin-provided stores that may evolve across
releases. Keep stable keys stable. If the durable store is the same, preserve
the key and update schema metadata; changing the key declares a new allocation
identity.

The normal integration pattern is to declare ranges with
\texttt{ic\_memory\_range!}, open through \texttt{ic\_memory\_key!}, and invoke
\texttt{bootstrap\_default\_memory\_manager()} from \texttt{init} and
\texttt{post\_upgrade} before any stable structure is touched.

Exactly one layer should bootstrap a given ledger store. Framework stacks should
compose declarations into that owner, or use distinct ledger stores and
allocation domains. Omitting a historical declaration does not retire or free
its key or slot; explicit retirement creates a tombstone and still does not make
the slot reusable for a different stable key.

\section{Limitations and Non-Goals}

The Lean model proves abstract allocation-ownership lemmas, and the Rust crate
enforces the corresponding invariants through recovery, validation, staging,
commit, and tests. This is not a formal verification of the Rust implementation.
It also does not prove semantic correctness of stored bytes, provide
cryptographic tamper resistance, protect against malicious controllers, validate
schema migrations, authorize endpoints, or provide disaster recovery. Its
checksum protects against torn writes and accidental corruption only.

The present native slot model follows \texttt{ic-stable-structures}
\texttt{MemoryManager} IDs exactly. Moving beyond 255 virtual memories would
require a different slot descriptor and memory-manager layout, not merely a
different \texttt{ic-memory} policy.

\section{Conclusion}

\texttt{ic-memory} turns stable-memory allocation from implicit convention into
durable protocol state. Its central guarantee is simple: a stable key cannot
move, and a physical slot cannot be silently reused. By enforcing that guarantee
before memory handles are opened, the crate closes a practical upgrade hazard
for multi-store Internet Computer canisters.

\end{multicols}

\end{document}
